\documentclass[aspectratio=169]{beamer}
\usetheme{SimplePlus}

\usepackage{tikz}
\usepackage{amsmath}
\usepackage{amssymb}
%\usepackage{bm}

\title[Finite Differences: Time Discretization]{Time Discretization and Stability}
\author{Nicol\'as A. Barnafi}
\date{}

\begin{document}

\begin{frame}
    \titlepage
\end{frame}

\begin{frame}{Outline}
    \tableofcontents
\end{frame}

\section{Time Discretization as a Quadrature Rule}

\begin{frame}{The Continuous and Semi-Discrete Problem}
    Consider a differential operator $A$ and its associated discrete matrix $A_h$.
    
    The continuous differential equation for an initial condition $x(0) = x_0$ is:
    \[
        \dot{x} + Ax = f
    \]
    Discretizing only in space yields a system of ODEs:
    \[
        \dot{x}_h + A_h x_h = f_h
    \]
\end{frame}

\begin{frame}{Time Discretization via Quadrature}
    Integrate the semi-discrete equation over  $(t_n, t_{n+1})$.  Using $x(t_n) \approx x^n$, this yields an exact relation:
    \[
        x^{n+1} - x^n + \int_{t_n}^{t_{n+1}} A_h x_h(s) ds = \int_{t_n}^{t_{n+1}} f_h(s) ds
    \]
    The choice of time-stepping scheme is entirely determined by how we approximate 
    $$\int_{t_n}^{t_{n+1}} A_h x_h(s) ds$$ 
    Euler, Runge-Kutta, etc.
\end{frame}

\begin{frame}{Three Fundamental Schemes}
    \begin{itemize}
        \item \textbf{Left-sided rule (Explicit):} Evaluates the operator at the known time $t_n$.
        \[
            \int_{t_n}^{t_{n+1}} A_h x_h(s) ds \approx \Delta t A_h x_h^n
        \]
        \item \textbf{Right-sided rule (Implicit):} Evaluates the operator at the unknown time $t_{n+1}$.
        \[
            \int_{t_n}^{t_{n+1}} A_h x_h(s) ds \approx \Delta t A_h x_h^{n+1}
        \]
        \item \textbf{Trapezoidal rule (Mid-point / Crank-Nicolson):} Averages the endpoints for higher accuracy.
        \[
            \int_{t_n}^{t_{n+1}} A_h x_h(s) ds \approx \frac{\Delta t}{2} (A_h x_h^{n+1} + A_h x_h^n)
        \]
    \end{itemize}
\end{frame}

\begin{frame}{The $\theta$-Method Generalization}
    A natural generalization of these three schemes is the $\theta$-method. For a fixed parameter $\theta \in [0, 1]$:
    \[
        \int A_h x_h(s) ds \approx \theta \Delta t A_h x_h^{n+1} + (1 - \theta) \Delta t A_h x_h^n 
    \]
    \vspace{0.3cm}
    \begin{itemize}
        \item $\theta = 0$ corresponds to the explicit scheme.
        \item $\theta = 1$ corresponds to the implicit scheme.
        \item $\theta = 1/2$ corresponds to the mid-point (Crank-Nicolson) scheme.
    \end{itemize}
\end{frame}

\section{Time Stability and the Lax Equivalence}

\begin{frame}{Time Stability Definition}
    Consider a general forward scheme written as $u^{n+1} = g(u^n)$.
    
    A condition such as $\|u^{n+1}\| \le c\|u^n\|$ for $c \le 1$ implies recursively that:
    \[
        \|u_h^n\| \le \|u_h^0\|
    \]
    This is precisely a stability result as in the \textbf{Lax Equivalence Theorem}.
\end{frame}

\begin{frame}{Stability of the First-Order Model Equation}
    Since our theory works for linear operators, we analyze stability using a simple scalar first-order equation (with $f \equiv 0$ for stability purposes):
    \[
        u' + cu = 0, \quad u(0) = u_0 
    \]
    \textbf{1. Explicit Scheme:}
    \[
        u^{n+1} = (1 - c\Delta t)u^n \implies |u^{n+1}| \le |1 - c\Delta t| |u^n| 
    \]
    The scheme is conditionally stable ($|1 - c\Delta t| < 1$) if $c > 0$ and $\Delta t < \frac{2}{c}$.
    
    \textbf{2. Implicit Scheme:}
    \[
        (1 + c\Delta t)u^{n+1} = u^n \implies |u^{n+1}| \le \left| \frac{1}{1 + c\Delta t} \right| |u^n| 
    \]
    If $c > 0$, then $1 + c\Delta t > 1$, meaning the method is unconditionally stable.
\end{frame}

\begin{frame}{Stability of the $\theta$-Method}
    Applying the $\theta$-method to the model problem yields:
    \[
        (1 + \theta c \Delta t)u^{n+1} = (1 - (1-\theta)c\Delta t)u^n 
    \]
    This gives the amplification bound:
    \[
        |u^{n+1}| \le \left| \frac{1 - (1-\theta)c\Delta t}{1 + \theta c\Delta t} \right| |u^n| 
    \]
    \textbf{Conclusions:}
    \begin{itemize}
        \item If $\theta \ge 1/2$, the fraction is bounded by 1, making the method \textbf{unconditionally stable}.
        \item If $\theta < 1/2$, stability requires a small $\Delta t$, meaning it is \textbf{conditionally stable}.
    \end{itemize}
\end{frame}

\section{Von Neumann Stability Analysis}

\begin{frame}{Moving to the Frequency Domain}
    Energy methods for stability analysis can quickly become intractable. To obtain easier bounds, we analyze the scheme in the frequency domain using \textbf{Von Neumann stability analysis}.
    
    This relies on the Fourier transform, utilizing the property that translation-invariant finite difference operators have simple eigenfunctions. 
    
    We consider a basis function for a given wave number $\xi$ evaluated on the grid $x_j = jh$:
    \[
        W_j = e^{i\xi(jh)} 
    \]
    This ansatz decouples the degrees of freedom, allowing us to find an amplification factor $g(\xi)$ such that $\hat{u}^{n+1} = g(\xi)\hat{u}^n$.
\end{frame}

\begin{frame}{Application: 1D Heat Equation}
    Consider the 1D heat equation $\dot{u} - u_{xx} = 0$ discretized with an explicit forward scheme in time and centered differences in space:
    \[
        \frac{u_j^{n+1} - u_j^n}{\Delta t} + \frac{1}{h^2}(-u_{j+1}^n + 2u_j^n - u_{j-1}^n) = 0 
    \]
    We introduce the ansatz $u_j^{n+1} = g(\xi) u_j^n = g(\xi) e^{i\xi(jh)}$. Substituting this into the scheme:
    \[
        g(\xi) e^{i\xi jh} = e^{i\xi jh} - \frac{\Delta t}{h^2} \left( -e^{i\xi(j-1)h} + 2e^{i\xi jh} - e^{i\xi(j+1)h} \right) 
    \]
\end{frame}

\begin{frame}{Extracting the Amplification Factor}
    Dividing the entire equation by $e^{i\xi j h}$, we isolate $g(\xi)$:
    \begin{align*}
        g(\xi) &= 1 + \frac{\Delta t}{h^2}(e^{-i\xi h} + e^{i\xi h} - 2)  \\
               &= 1 + \frac{\Delta t}{h^2}(2\cos(\xi h) - 2) \quad \text{(using Euler's formula)}  \\
               &= 1 - \frac{2\Delta t}{h^2}(1 - \cos(\xi h)) 
    \end{align*}
    For the scheme to be stable, we require $|g(\xi)| \le 1$ for all wave numbers $\xi$.
\end{frame}

\begin{frame}{Deriving the Stability Condition (CFL)}
    Since $\cos(\xi h) \in [-1, 1]$, the term $(1 - \cos(\xi h))$ ranges from $0$ to $2$. This yields two bounds for $g(\xi)$:
    
    \vspace{0.3cm}
    \textbf{Upper bound (always satisfied):}
    \[
        g(\xi) \le 1 - \frac{2\Delta t}{h^2}(0) = 1 
    \]
    
    \textbf{Lower bound (imposes a restriction):}
    \[
        g(\xi) \ge 1 - \frac{2\Delta t}{h^2}(2) = 1 - \frac{4\Delta t}{h^2} 
    \]
    For stability, we must enforce $-1 \le g(\xi)$.
\end{frame}

\begin{frame}{Deriving the Stability Condition (CFL)}
    Solving the inequality for the lower bound:
    \begin{align*}
        1 - \frac{4\Delta t}{h^2} &\ge -1  \\
        2 &\ge \frac{4\Delta t}{h^2}  \\
        \Delta t &\le \frac{h^2}{2} 
    \end{align*}
    This final inequality is the stability threshold for the explicit finite difference scheme of the heat equation. If the time step exceeds this upper bound, the amplification factor grows beyond $1$, and the numerical solution becomes unconditionally unstable.
\end{frame}

\end{document}

\begin{frame}
    \maketitle
\end{frame}
\end{document}
